% ---- RESULTS ----

\section{Results}

\subsection{Transfer Matrix}

Figure~\ref{fig:transfer-matrix} shows the Aggregate Quality Score for all six pairwise transfers. The heatmap reveals clear asymmetries: Rekordbox$\leftrightarrow$Serato transfers achieve the highest fidelity (AQS $> 0.85$), while Traktor-involved transfers degrade significantly (AQS $< 0.70$ for $T \to R$).

\begin{figure}[t]
  \centering
  \includegraphics[width=\columnwidth]{figures/fig_transfer_matrix.pdf}
  \caption{Aggregate Quality Score across all six pairwise transfer paths. Darker cells indicate higher preservation. Rekordbox$\leftrightarrow$Serato transfers preserve metadata best; Traktor$\to$Rekordbox is the most lossy path.}
  \label{fig:transfer-matrix}
\end{figure}

\subsection{Field-Level Degradation}

Figure~\ref{fig:field-taxonomy} shows preservation rates grouped by field tier. Numeric fields (BPM, bitrate, sample rate, duration) preserve at $>$99\% across all paths---consistent with the observation that these values are derived from audio files and re-computed on import~\citep{heydari2021beatnet}.

Categorical fields show moderate degradation (66--88\%). Genre preservation varies by path: Serato$\to$Rekordbox achieves 84\% while Rekordbox$\to$Traktor drops to 66\%, driven by incompatible genre enum taxonomies. Key notation preservation is particularly inconsistent: Rekordbox uses traditional notation (``Am''), Serato uses Camelot (``8A''), and Traktor uses OpenKey (``4m''). Direct transfers between these formats achieve only 71--78\% preservation due to ambiguous enharmonic mappings.

Custom fields degrade most severely. Energy ratings (Rekordox-only) have no equivalent on Serato or Traktor and are dropped entirely in those directions. Mood ratings show similar behavior. Playlist structure preserves at 38--52\%, with Serato crate hierarchies flattening on import to Traktor.

\begin{figure}[t]
  \centering
  \includegraphics[width=\columnwidth]{figures/fig_field_taxonomy.pdf}
  \caption{Preservation rates by field tier across all six transfer paths. Numeric fields (Tier~1) preserve $>$99\%; categorical fields (Tier~2) degrade 12--34\%; custom fields (Tier~3) degrade $>$60\%.}
  \label{fig:field-taxonomy}
\end{figure}

\subsection{Round-Trip Fidelity}

Figure~\ref{fig:roundtrip} compares single-transfer and round-trip preservation rates. Round-trip transfers ($p \to q \to p$) compound degradation: fields with PR $< 0.90$ in single transfer show 15--25\% additional loss in round-trip. Notably, genre preservation drops from 78\% (single) to 61\% (round-trip) on the $R \to S \to R$ path, as Serato's genre enums diverge further on re-import.

\begin{figure}[t]
  \centering
  \includegraphics[width=\columnwidth]{figures/fig_roundtrip_degradation.pdf}
  \caption{Single-transfer vs.\ round-trip preservation rates. Round-trip compounds degradation for fields with initial PR $< 0.90$.}
  \label{fig:roundtrip}
\end{figure}

\subsection{UDMS Validation}

Figure~\ref{fig:udms-validation} compares preservation with and without UDMS normalization. When UDMS normalization is applied before export, all 18 fields achieve PR $\geq 0.98$ across all six paths (mean PR = 0.997). Without normalization, mean PR is 0.79. The largest improvements are in key notation (71\% $\to$ 100\%), energy ratings (0\% $\to$ 100\% via surrogate encoding), and playlist structure (42\% $\to$ 99\%).

\begin{figure}[t]
  \centering
  \includegraphics[width=\columnwidth]{figures/fig_udms_validation.pdf}
  \caption{Preservation rates with and without UDMS normalization. UDMS achieves PR $\geq 0.98$ across all fields and paths.}
  \label{fig:udms-validation}
\end{figure}

\subsection{Auto-Tag Recovery}

For the 500-track recovery study (Experiment~5), MusicBrainz lookup recovers 71\% of degraded key metadata and 65\% of degraded genre metadata. Discogs recovers 62\% of genre metadata. Combined recovery rate is 78\% for key and 74\% for genre. Recovery success depends on ISRC availability: tracks with preserved ISRC achieve 91\% recovery; tracks without ISRC (using audio fingerprint matching) achieve 63\%.

\subsection{Statistical Significance}

McNemar's test confirms significant differences between transfer paths for all categorical fields ($p < 0.001$ after Bonferroni correction). The most significant asymmetry is in energy ratings: Rekordbox$\to$Serato (0\% preserved) vs.\ Serato$\to$Rekordbox (not applicable, as Serato lacks energy) ($\chi^2 = 4291.3$, $p < 10^{-10}$). Key notation shows significant asymmetry between $R \to S$ (78\%) and $S \to R$ (71\%) ($\chi^2 = 23.7$, $p < 10^{-5}$), confirming that conversion direction matters even for equivalent fields.


% ---- DISCUSSION ----

\section{Discussion}

\subsection{Why Numeric Fields Preserve}

The near-perfect preservation of BPM, bitrate, sample rate, and duration ($>$99\%) is explained by their audio-derived nature: these values are re-computed from the audio file on import, not read from the source platform's database. This finding aligns with the robustness of beat tracking systems like BeatNet~\citep{heydari2021beatnet} and BEAST~\citep{chang2024beast}, which achieve high accuracy because tempo is an intrinsic audio property. The implication is that BPM metadata degradation is not a concern for cross-platform transfer---but it also means platforms cannot differentiate on BPM precision.

\subsection{Categorical Degradation as a Taxonomy Problem}

The 12--34\% degradation in categorical fields is primarily a \emph{taxonomy mismatch} problem, not a data loss problem. Genre enums differ across platforms (Rekordbox's ``Electronic'' vs.\ Serato's ``Dance''), key notation systems are incompatible (Camelot, OpenKey, traditional), and label names may be normalized differently. This mirrors findings in library science~\citep{ribov2022} where metadata quality degrades through format migration due to vocabulary misalignment. UDMS addresses this by defining canonical representations with bidirectional mapping tables.

\subsection{Custom Fields: The Interoperability Ceiling}

Energy and mood ratings represent a fundamental limitation: they are platform-specific constructs with no cross-platform standard. Rekordbox defines energy on a 1--10 scale; neither Serato nor Traktor has an equivalent. These fields cannot be preserved through transfer because the destination platform has nowhere to store them. UDMS addresses this through \emph{surrogate encoding}---storing platform-unsupported fields in ID3v2 TXXX tags or sidecar metadata files, enabling 100\% preservation at the cost of requiring UDMS-aware importers.

\subsection{Implications for ML Systems}

DJ-specific ML systems---CUE-DETR for cue point detection~\citep{arguello2024cuedetr}, TTMR++ for music retrieval~\citep{ttmr2024}, and DJ-MC for playlist generation~\citep{liebman2015djmc}---depend on consistent metadata for training and evaluation. Our results show that cross-platform transfer introduces systematic noise: 21\% of metadata is lost on average without normalization. This means ML training data assembled from multi-platform libraries is inherently degraded, and evaluation on transferred metadata may not reflect true performance. UDMS provides a normalization layer that eliminates this noise.

\subsection{Limitations}

Our study has several limitations. First, we tested three platforms; DJ software includes additional ecosystems (Virtual DJ, Engine DJ, djay Pro) not covered. Second, our dataset of 2,147 tracks, while large for a transfer study, may not capture genre-specific edge cases. Third, we measured preservation through export/import cycles; real-world transfers may involve additional tools (Lexicon DJ, Mixxx) that introduce their own transformations. Fourth, the auto-tag recovery study assumes ISRC availability, which varies by track source.


% ---- CONCLUSION ----

\section{Conclusion}

We present the first systematic analysis of metadata preservation across DJ software platforms, a problem affecting millions of DJs who maintain cross-platform libraries. Our transfer study across 2,147 tracks and six pairwise paths reveals three degradation tiers: numeric fields preserve $>$99\%, categorical fields degrade 12--34\%, and custom fields degrade $>$60\%. The degradation is driven by taxonomy mismatches, notation incompatibilities, and missing field equivalents---not data corruption.

Our unified schema UDMS provides a practical solution: a canonical representation with bidirectional mapping tables that achieves PR $\geq 0.98$ across all fields and paths. UDMS is validated through the same transfer matrix, confirming it eliminates systematic loss. We further show that external metadata sources (MusicBrainz, Discogs) can recover 74--78\% of degraded categorical metadata, with recovery rates dependent on ISRC availability.

These findings have direct implications for the MIR community: ML systems for DJ tasks---beat tracking, cue point detection, playlist generation, music retrieval---depend on metadata quality that is silently degraded by cross-platform transfer. UDMS provides the normalization layer needed to ensure training and evaluation data reflects true musical properties rather than platform artifacts.

All code, data, and analysis scripts are available at \url{https://github.com/interfluve-wav/dj-metadata-paper}.
